\section{Introduction}

A \((k,n)\)-arc in a projective plane is a set of \(k\) points meeting every
line in at most \(n\) points, with some line meeting it in exactly \(n\)
points.  It is complete if it is maximal under inclusion, or equivalently if
every external point lies on an \(n\)-secant.  The minimum size of a complete
\((k,n)\)-arc is denoted by \(t_n(2,q)\) in
\(\PG(2,q)\).  Lower bounds for this quantity are usually obtained by counting
maximal secants or, dually, by bounding minimal multiple blocking sets; see
Alabdullah--Hirschfeld~\cite{AlabdullahHirschfeld} and
Bishnoi--Mattheus--Schillewaert~\cite{BMS}.  Recent work has also developed
algebraic and curve-based methods for proving completeness itself; see, for
example, Korchm\'aros--Nagy--Sz\H{o}nyi~\cite{KNS}.

This paper keeps the maximal secants as a family rather than eliminating their
number at the first counting step.  The degrees of that family on the arc and
on its complement have fixed sums, and their pair counts add to the number of
pairs of secants.  Integer convexity gives the exact range of each pair count.
The two ranges must overlap.  Replacing the integer extrema by quadratic real
averages recovers the classical incidence-mixing inequality; retaining them
gives the new bounds.

The strongest application concerns ordinary completeness in characteristic
three.

\begin{theorem}[Characteristic-three bound]\label{thm:characteristic-three}
\coverage{absent}
For every \(\varepsilon>0\), there is \(h_0\) such that, for \(h\ge h_0\),
every complete \((k,2q/3+1)\)-arc in \(\PG(2,q)\), where \(q=3^h\), satisfies
\begin{equation}\label{eq:characteristic-three-bound}
 k\ge \frac{q^2}{3}+\left(\frac43-\varepsilon\right)q.
\end{equation}
Equivalently,
\[
 t_{2q/3+1}(2,q)\ge \frac{q^2}{3}+\frac{4q}{3}-o(q)
 \qquad (q=3^h\to\infty).
\]
\imports{SWmodp:repair}
\end{theorem}

The incidence-mixing bound gives only \(q^2/3+q+O(1)\) in this regime.  The
integer degree bounds first raise the linear coefficient to \(6/5\); a
modular repair theorem of Sz\H{o}nyi--Weiner~\cite{SWmodp} then raises it to
\(4/3\).  The second step uses more than modular stability: each point changed
in the dual repaired multiset becomes a primal line on which completeness
forces a linear amount of external secant excess.

The arithmetic source of the first improvement is visible whenever the real
incidence bound is tight at a rational parameter.  Fix a coverage multiplicity
\(\lambda\), factor it as
\(\lambda=uv\), and put \(d=u+v+1\).  For \(h\in\mathbb Z\), define
\begin{align}
 c_{\mathrm{bal}}
 &=\frac{u(uv^2+4uv+2u+2v^2+4v+1)}{2uv+u+v},\label{eq:c-balanced}\\
 h_*&=-\frac{u(v+1)(u^2v+u^2-v)}{2uv+u+v},\label{eq:h-star}\\
 \mu&=\frac{2v}{(u+v)(2uv+u+v+1)},\nonumber\\
 L(h)&=d-\frac hu,\qquad
 C(h)=c_{\mathrm{bal}}+\mu(h-h_*),\nonumber\\
 c_{\mathbb Z}&=\min_{h\in\mathbb Z}\max\{L(h),C(h)\}.
 \label{eq:c-integer}
\end{align}

\begin{theorem}[Rational equality families]\label{thm:factor-pair}
\coverage{absent}
Let \(u,v\ge1\), \(\lambda=uv\), and \(d=u+v+1\).  Suppose that a
\((k,(u+1)m+1)\)-arc in a projective plane of order \(q=dm\) has at least
\(\lambda\) \(n\)-secants through every external point.  Then
\begin{equation}\label{eq:factor-pair-bound}
 k\ge udm^2+c_{\mathbb Z}m-O_{u,v}(1).
\end{equation}
Moreover, the classical incidence-mixing bound has linear coefficient
\[
 c_{\mathrm{mix}}=\frac{u(uv+3u+2v+3)}{2u+1},
\]
and
\begin{equation}\label{eq:strict-gap}
 c_{\mathrm{bal}}-c_{\mathrm{mix}}
 =\frac{u(v+1)(u^2+u+1)}{(2u+1)(2uv+u+v)}>0.
\end{equation}
Every rational equality parameter of the continuous bound arises from exactly one
ordered factorization \(\lambda=uv\).
\imports{BMS:spectral-bound}
\end{theorem}

Thus the strict improvement is not an isolated rounding effect.  At fixed
coverage multiplicity there is one equality family for each divisor, and every
family has a positive linear integrality penalty.  The two integers adjacent
to \(h_*\) determine \(c_{\mathbb Z}\), so the coefficient is explicit.

The same integer incidence condition holds in arbitrary symmetric designs.
In a projective plane there are two further bounds.  At most
\(\lfloor(k-1)/(n-1)\rfloor\) \(n\)-secants pass through a point of the arc,
and at most \(\lfloor k/n\rfloor\) pass through an external point.  These
geometric bounds sharpen the general condition for complete arcs.

In characteristic two the first family in which the two first-order bounds
are simultaneously tight yields a second
application.

\begin{theorem}[Characteristic-two parity bound]\label{thm:characteristic-two}
\coverage{absent}
For every \(\varepsilon>0\), there is \(h_0\) such that, for \(h\ge h_0\),
if \(q=2^h\), \(8\mid q\), and every external point of a
\((k,3q/4+1)\)-arc in \(\PG(2,q)\) lies on at least two \((3q/4+1)\)-secants,
then
\begin{equation}\label{eq:characteristic-two-bound}
 k\ge \frac{q^2}{2}+\left(\frac{73}{48}-\varepsilon\right)q.
\end{equation}
\imports{SWeven:repair}
\end{theorem}

The proof separates odd and even repair support.  The parity of the selected
secant family fixes the parity of the support, while the symmetric difference
of its generator lines supplies respectively \(q/4-O(1)\) or \(q/2-O(1)\)
external points of excess degree.

The projective-code translation is immediate.  A point set in
\(\PG(r-1,q)\) gives a projective \([k,r,k-s]_q\) code; maximal hyperplanes
are minimum-weight codewords, and a candidate extension column is obstructed
by each maximal hyperplane through it.  The same integer degree theorem
therefore bounds robust nonextendibility.  We keep the finite-geometry
language primary and record this translation after the planar applications.

The paper is organized as follows.  Section~\ref{sec:integer-envelopes}
proves the symmetric-design inequality and identifies incidence mixing as its
real relaxation.  The next two sections specialize to maximal secants and
prove Theorem~\ref{thm:factor-pair}.  The modular-lift section proves
Theorems~\ref{thm:characteristic-three} and
\ref{thm:characteristic-two}.  The final section gives the blocking-set and
coding translations and records the exact computational checks used to guard
the finite algebra.
